% =============================================================================
% Section 3: Methods
% =============================================================================

\section{Methods}
\label{sec:methods}

\subsection{Model Architecture}
\label{sec:arch}

All models follow a standard decoder-only transformer architecture \cite{vaswani2017attention,radford2019language}
with the following fixed design choices: pre-layer normalization \cite{xiong2020layer}, weight-tied
input and output embeddings, no bias terms in any linear projection, no dropout, and a learned positional
embedding table. These choices simplify the residual stream algebra used in circuit analysis
\cite{elhage2021mathematical}.

The configurable axes are depth $L$ (number of blocks), width $d_{\text{model}}$, feed-forward
expansion $d_{\text{ff}}$, and number of attention heads $H$. Each attention block implements
$H$ heads in parallel. For head $h$, the attention-weighted value is

\begin{equation}
\label{eq:attn}
\text{attn}_h(x) = \operatorname{softmax}\!\left(\frac{Q_h K_h^T}{\sqrt{d_k}}\right) V_h,
\end{equation}

\noindent where $Q_h = x W_Q^{(h)}$, $K_h = x W_K^{(h)}$, $V_h = x W_V^{(h)}$, and
$d_k = d_{\text{model}} / H$. A causal mask is applied before the softmax to enforce
autoregressive conditioning. All heads are projected and summed via a shared output matrix $W_O$.

Attention weights at each layer are cached during inference to support the circuit-extraction
methods described in Section~\ref{sec:circuit_methods}. Table~\ref{tab:model_configs} lists
the specific configurations used in each experiment.

\begin{table}[h]
\centering
\caption{Model configurations used across experiments.}
\label{tab:model_configs}
\begin{tabular}{lrrrrrr}
\toprule
Experiment & $L$ & $H$ & $d_{\text{model}}$ & $d_{\text{ff}}$ & Vocab & Params \\
\midrule
Positional induction (seed divergence) & 4 & 4 & 128 & 512 & 50 & $\approx$800K \\
IOI (seed divergence) & 4 & 4 & 128 & 512 & variable & $\approx$800K \\
Content-matching induction & 2 & 4 & 128 & 512 & 512 & $\approx$460K \\
Markov entropy continuum & 4 & 4 & 128 & 512 & 64 & $\approx$800K \\
Shakespeare (6-layer) & 6 & 8 & 512 & 2048 & 65 & $\approx$19M \\
Depth scaling ($L \in \{2,4,6,8\}$) & var & 8 & 512 & 2048 & 65 & 6.3M--25.2M \\
\bottomrule
\end{tabular}
\end{table}

\subsection{Training Protocol}
\label{sec:training}

All models are trained with AdamW \cite{loshchilov2018decoupled} using a cosine learning rate
schedule with linear warmup. Gradient entries are clipped elementwise to $[-1, 1]$ before
each parameter update. The default learning rate is $10^{-3}$ with weight decay in $[0.01, 0.1]$
(see Appendix~\ref{app:training_details} for per-experiment values). There is no dropout, as our
primary interest is the converged weight structure rather than generalization.

The loss function is masked cross-entropy: on tasks with sparse prediction targets (induction,
IOI), we mask the loss to positions where the model must predict a specific non-trivial token,
ignoring filler positions where any prediction is acceptable. On dense prediction tasks
(Shakespeare, Markov chains), loss is computed over all positions.

Checkpoints are saved every 25 training steps for the seed-divergence experiments and every
100--500 steps for all other experiments, enabling post-hoc analysis of circuit formation
dynamics (Section~\ref{sec:formation}). All training is performed using the MLX framework
\cite{mlx2023} on Apple Silicon (M-series) hardware using the Metal GPU backend.
The total compute budget for all reported experiments is approximately six hours of wall-clock time.

\subsection{Tasks}
\label{sec:tasks}

We study five task families chosen to span a range of computational demands.

\paragraph{Positional induction.}
Sequences have the form $[T_1, \ldots, T_n, T_1, \ldots, T_n]$ where $T_i$ are drawn uniformly
from a vocabulary of size 50 and $n = 32$ (total sequence length 64). Every position in the
second half is a prediction target: the model must output the token at the corresponding position
in the first half. This implements the Olsson et al.\ \cite{olsson2022context} dense induction
signal and avoids the positional-sparsity failure mode. Because the repeated token always appears
at a fixed positional offset ($i - n$), a pure positional attention mechanism is sufficient;
no content matching is required.

\paragraph{Content-matching induction.}
To force the model to develop genuine token-identity-based induction, we use random-position
bigrams. Each sequence of length 64 contains $k = 6$ token pairs $(A_i, B_i)$ placed at two
randomly chosen positions each. The model must predict $B_i$ when it encounters the second
occurrence of $A_i$. Vocabulary size is 512, so random token collisions are low but nonzero
under the fixed generator and evaluation seeds. No fixed positional shortcut exists; the model
must attend to the first occurrence of the matching token.

\paragraph{Indirect object identification (IOI).}
Synthetic sentences of the form \textit{``[Name1] and [Name2] went to [place]. [Name2] gave the
[object] to''} require the model to predict \textit{Name1} (the indirect object). We use a
pool of 20 first names, 10 action verbs, and 30 filler words, generating sequences programmatically.
This task requires suppression of the subject name (Name2) and promotion of the indirect object
(Name1), matching the structure of the GPT-2 IOI circuit \cite{wang2022interpretability}.

\paragraph{Markov chains with varying entropy.}
A first-order Markov chain over vocabulary size 64 is parameterized by an entropy level
$\alpha \in [0, 1]$. At $\alpha = 0$ the chain is deterministic (each token has exactly one
successor); at $\alpha = 1$ the chain is uniform (each token is equally likely to follow any
other). Intermediate values interpolate via temperature scaling of the transition matrix.
We train 11 models at $\alpha \in \{0.0, 0.1, 0.2, \ldots, 1.0\}$ with sequences of length 64.
This family allows us to measure how circuit architecture varies continuously with task entropy.

\paragraph{Shakespeare (natural language).}
Character-level language modeling on the TinyShakespeare corpus \cite{karpathy2015char},
consisting of approximately $1.1 \times 10^6$ characters from Shakespeare's plays.
Vocabulary size is 65 (printable ASCII characters). This task serves as an ecological validity
check: do circuit motifs identified on synthetic tasks generalize to natural language?

\subsection{Circuit Extraction Methods}
\label{sec:circuit_methods}

We use four complementary tools to map and quantify circuits.

\paragraph{Activation patching.}
For a clean input $x_{\text{clean}}$ and a corrupted input $x_{\text{corr}}$, activation
patching measures the causal contribution of component $(i, j)$ (layer $i$, component type $j$)
by replacing its activations with the corrupted version and measuring the logit difference:
\begin{equation}
\label{eq:patching}
\Delta_{ij} = \ell_{\text{clean}} - \ell_{\text{patched}(i,j)},
\end{equation}
where $\ell$ denotes the logit of the target token. A large $\Delta_{ij}$ indicates the component
is causally load-bearing for the prediction.

\paragraph{Direct logit attribution.}
Each component $c$'s contribution to the final logit for target token $t$ is computed as
\begin{equation}
\label{eq:logit_attr}
\text{attr}_c = \langle \text{contrib}_c,\, W_U[t] \rangle,
\end{equation}
where $\text{contrib}_c$ is the component's output projected to the residual stream basis and
$W_U[t]$ is the unembedding vector for token $t$. Because the architecture uses weight tying,
$W_U = W_E^T$. Positive attribution means the component promotes the target token; negative
attribution means it suppresses competitors or the target itself.

\paragraph{OV circuit decomposition.}
The output-value (OV) matrix of head $h$ captures what the head copies or transforms when it
attends. Written in terms of the full embedding space:
\begin{equation}
\label{eq:ov}
\text{OV}_h = W_E\, W_V^{(h)}\, W_O^{(h)}\, W_E^T.
\end{equation}
A head implements identity copying if this matrix is close to diagonal: the token attended to
is the token written to the residual stream. We measure the \emph{copy score} as the mean diagonal
of $\text{OV}_h$ minus the mean off-diagonal, normalized by the off-diagonal standard deviation.
A high positive copy score indicates the head aggressively copies the token identity it attends
to; a negative score (suppression) indicates the head promotes tokens \emph{other} than the one
it attends to.

\paragraph{Permutation-sensitivity metric.}
Standard positional-dependence metrics measure how much a head's attention pattern varies across
different inputs at fixed positions. This metric correctly identifies positional heads (which attend
to fixed positions regardless of input content) but systematically misclassifies content-based heads
in deep layers where the residual stream already carries rich, stable semantic structure. The
attention pattern can be highly correlated across inputs while still being content-conditioned, if
the content is itself stable.

We introduce a permutation-sensitivity metric that directly tests whether a head's attention pattern
depends on token identity. Given a batch of inputs $X$, we apply a random vocabulary permutation
$\pi$ (shuffling the token-to-embedding mapping) and measure how much the head's attention
matrix changes:
\begin{equation}
\label{eq:perm_sens}
\text{sens}_h = \frac{1}{N} \sum_{\pi} \|A_h(X) - A_h(\pi(X))\|_F,
\end{equation}
where $A_h(\cdot)$ denotes the attention weight matrix of head $h$ and the average is over $N$
random permutations. A content-based head will show large sensitivity (its pattern changes when
token identities are shuffled); a truly positional head will show near-zero sensitivity (it
attends to fixed positions regardless of token content). This metric resolves the misclassification
of deep-layer suppression heads that the standard $\text{pos\_dep}$ score misidentifies as
positional (Section~\ref{sec:metric_correction}; Appendix~\ref{app:metric_validation}).

\subsection{Surgical Ablation Protocol}
\label{sec:ablation}

To causally validate circuit maps, we zero out individual components during the forward pass.
For a single attention head $(l, h)$, we multiply the head's value-weighted output by zero
\emph{before} the output projection:
\begin{equation*}
\tilde{o}_{l,h} = 0 \cdot (A_{l,h}\, V_{l,h}),
\end{equation*}
leaving all other heads' contributions to the output projection unchanged. For whole-layer
ablations (all heads in layer $l$, or the MLP in layer $l$), the corresponding residual stream
contribution is zeroed entirely.

Because MLX \cite{mlx2023} uses lazy, functional evaluation and does not support in-place tensor
mutation, ablations are implemented by constructing a head-selection mask at run time:
\texttt{head\_mask} $= [0 \text{ if } h = h_{\text{ablate}} \text{ else } 1]$, reshaped to
broadcast over the \texttt{(B, H, T, d\_head)} output tensor. Multiple simultaneous ablations are
supported by iterating over the ablation set before the output projection. Full pseudocode is given
in Appendix~\ref{app:ablation_protocols}.

We report task accuracy (for induction tasks) or cross-entropy loss (for Shakespeare) relative to
the unablated baseline. A component is considered \emph{individually critical} if its ablation
causes an accuracy drop exceeding 0.5 (i.e., more than halving the task accuracy).

\subsection{Transplantation Protocol}
\label{sec:transplant}

To test whether circuit weights are portable across independently trained networks, we implement
a surgical weight-transplantation procedure. Given a donor network (seed 0) with critical
head at $(l_0, h_0)$ and a recipient network (seed 1) with critical head at $(l_1, h_1)$, we:
\begin{enumerate}
    \item Extract the weight slices $W_Q^{(h_0)}, W_K^{(h_0)}, W_V^{(h_0)}$ (rows in the fused
          QKV projection) and the corresponding columns of $W_O^{(l_0)}$ from the donor.
    \item Overwrite the corresponding slices in the recipient's $(l_1, h_1)$ slot.
    \item Verify the transplant via a parameter checksum: confirm the transplanted weights differ
          from the recipient's original weights and match the donor's.
\end{enumerate}

To test whether basis mismatch (rather than fundamental non-portability) explains transplant
failure, we additionally apply Procrustes basis alignment \cite{schonemann1966procrustes}.
We collect residual stream activations $X_0, X_1 \in \mathbb{R}^{n \times d}$ at the boundary
between layers $l-1$ and $l$ from both networks on identical inputs, then solve for the optimal
orthogonal rotation:
\begin{equation}
\label{eq:procrustes}
R^* = \arg\min_{R:\, R^T R = I} \|X_0 R - X_1\|_F.
\end{equation}
Via the SVD $X_1^T X_0 = U S V^T$, the solution is $R^* = U V^T$ \cite{horn1994topics}.
We apply $R^*$ to rotate the donor head's weight matrices into the recipient's activation basis
before transplantation, then re-evaluate task accuracy. A negative or negligible improvement
over the unaligned transplant indicates the two networks' residual stream bases are not related
by any orthogonal rotation.

\subsection{Statistical Methods}
\label{sec:stats}

All confidence intervals are bootstrap 95\% CIs with 1000 resamples. Agreement rates
(fraction of seed-pairs sharing the same critical head identity) are tested against a null of
$1/H$ (uniform random assignment over $H$ heads) using a two-sided permutation test with 10,000
permutations. Ablation drop means are tested against zero using a paired $t$-test across seeds.
The suppression-fraction scaling law is estimated by ordinary least squares:
$\text{frac\_suppression} \sim \beta_0 + \beta_1 \cdot L$, with bootstrap CIs on $\beta_1$.

Results are reported as point estimate [CI$_{\text{low}}$, CI$_{\text{high}}$] with $p$-value.
Statistical significance threshold is $\alpha = 0.05$ throughout. The four key statistics and
their bootstrap CIs are summarized in Table~\ref{tab:statistics}.

\begin{table}[h]
\centering
\caption{Key quantitative claims with bootstrap 95\% confidence intervals.}
\label{tab:statistics}
\begin{tabular}{lcccc}
\toprule
Claim & Estimate & 95\% CI & Null & $p$ \\
\midrule
Cross-seed circuit similarity & 0.985 & [0.982, 0.988] & 0 & $<0.001$ \\
Head-level agreement (20 seeds) & 13.2\% & [11.6\%, 27.4\%] & 6.25\% & 0.001 \\
Critical-head ablation drop (4 seeds) & 0.775 & [0.590, 0.906] & 0 & 0.004 \\
Suppression fraction slope (per layer) & 0.059 & [0.023, 0.109] & 0 & 0.052 \\
\bottomrule
\end{tabular}
\end{table}
